\documentclass[10pt,a4paper]{amsart}
\usepackage[margin=19mm]{geometry}
\usepackage{amsmath,amssymb,mathtools}
\usepackage[scr=rsfs,cal=euler]{mathalfa}
\usepackage{xfrac}
\usepackage[unicode,hidelinks,bookmarks=false]{hyperref}
\newcommand{\ie}{i.e.\ }
\newcommand{\cf}{cf.\ }
\newcommand{\resp}{respectively\ }
\newcommand{\Subsection}{\subsection}
\providecommand{\nobreakdash}{\nobreak\hbox{-}}
\setlength{\emergencystretch}{2em}
\multlinegap0pt
\usepackage{bm}
\usepackage{bbm}
\usepackage{dsfont}
\usepackage{xspace}
\usepackage{cancel}
\usepackage{mathtools}
\usepackage[shortlabels]{enumitem}
\usepackage{amsthm}
\makeatletter
\@ifundefined{theorem}{\newtheorem{theorem}{Theorem}}{}
\@ifundefined{lemma}{\newtheorem{lemma}[theorem]{Lemma}}{}
\makeatother
\DeclareMathOperator{\ord}{ord}
\title[Certified Minimal-Prime Branch Closures for Odd Perfect Numb]{Certified Minimal-Prime Branch Closures for Odd Perfect Numbers}
\author{Marco Mantovanelli}
\date{Source: arXiv:2607.04365v1 (CC BY 4.0)}
\begin{document}
\maketitle

\noindent\emph{Source and licence.} Certified Minimal-Prime Branch Closures for Odd Perfect Numbers. Version arXiv:2607.04365v1 (CC BY 4.0), distributed under the Creative Commons Attribution 4.0 International licence, as recorded in the arXiv licence metadata for this exact version and shown on the abstract page https://arxiv.org/abs/2607.04365v1. Full rights notice: licenses/arxiv-2607.04365-CC-BY-4.0.txt.\par\smallskip

\noindent\textit{Editorial scope.} Counted unit: the Lemma whose source label is \texttt{lem:q5-first-input-classification} in the article (source0.tex lines 1426--1503, source label \texttt{lem:q5-first-input-classification}), delivered together with the article's own complete original proof of it (source0.tex lines 1505--1623, one \texttt{proof} environment of 119 lines). No mathematics is written, repaired or re-derived: statement and proof are the article's own text line for line. Required context retained uncounted: none. Every definition, display and label of those lines is retained unchanged. Omitted from this package and not redistributed: 1--1425, 1504--1504, 1624--3798. Nothing inside a retained excerpt is omitted or altered: every retained body is delivered exactly as the article's lines are written, so no omission receipt of any kind accompanies this package. Citations: the retained excerpts carry no citation command at all, so no bibliography entry of the article is transcribed and no reference is redistributed. Reference adaptation: none was needed and none was made; every reference inside the delivered unit resolves inside it, so no unresolved reference or citation is produced. Printed slips kept literal: no printed word, symbol, hypothesis or number of the counted statement or of its proof was altered. Layout and class: the article's own class is \texttt{article} with packages including fontenc, inputenc, amsmath, amssymb, amsthm, mathtools, enumitem, geometry, booktabs, array, tabularx, placeins, hyperref; the wrapper is the lane's pinned \texttt{amsart} editorial wrapper at 10pt on A4, which restates the theorem-like environments the retained text needs and adds the article's own macro definitions that the retained text needs (\texttt{\textbackslash ord}), with extra wrapper packages (bm, bbm, dsfont, xspace, cancel, mathtools, enumitem). The delivered numbering is the unit's own. Assets: no retained span contains a figure environment, a graphics-inclusion command, a PDF-page inclusion, a table or a TikZ picture; no image file is copied into this package, the package declares no asset dependency and no image file is redistributed with it. Licence: version arXiv:2607.04365v1 of this article is distributed under the Creative Commons Attribution 4.0 International licence, as recorded in the arXiv licence metadata for this exact version and shown on the abstract page https://arxiv.org/abs/2607.04365v1. A full rights notice is delivered at licenses/arxiv-2607.04365-CC-BY-4.0.txt. Characters: every retained excerpt is pure ASCII, so the delivered engine, XeLaTeX with its default Latin Modern text font, reports no missing character.
\begin{lemma}[Classification of first $5$-input witnesses]
\label{lem:q5-first-input-classification}
Assume
\[
  \min S=5.
\]
Let
\[
  p\in S,\qquad p\ne5,
\]
be a $5$-input witness.

Then exactly one of the following role cases applies.

\begin{enumerate}[label=(\roman*),leftmargin=2.2em]

\item \textbf{Non-Euler witness.}
      The prime $p$ is non-Euler.  Then
      \[
        e_p\equiv0\pmod2,
      \]
      and the $5$-input condition is equivalent to
      \[
        p\equiv1\pmod5,
        \qquad
        5\mid e_p+1.
      \]
      In this case
      \[
        v_5(\sigma(p^{e_p}))=v_5(e_p+1).
      \]

\item \textbf{Euler witness of order one.}
      The prime $p$ is the Euler prime $\pi$, and
      \[
        \pi\equiv1\pmod5,
        \qquad
        5\mid \alpha+1.
      \]
      In this case
      \[
        v_5(\sigma(\pi^\alpha))=v_5(\alpha+1).
      \]

\item \textbf{Euler witness of order two.}
      The prime $p$ is the Euler prime $\pi$, and
      \[
        \pi\equiv-1\pmod5.
      \]
      In this case
      \[
        \ord_5(\pi)=2,
      \]
      and
      \[
        v_5(\sigma(\pi^\alpha))
        =
        v_5(\pi^2-1)
        +
        v_5\!\left(\frac{\alpha+1}{2}\right).
      \]
\end{enumerate}

Moreover, the Euler witness cannot have
\[
  \pi\equiv2\pmod5
  \qquad\text{or}\qquad
  \pi\equiv3\pmod5,
\]
because then
\[
  \ord_5(\pi)=4
\]
and
\[
  4\nmid \alpha+1.
\]
\end{lemma}

\begin{proof}
Let
\[
  n_p:=e_p+1.
\]

First suppose that $p$ is non-Euler.  Then Euler's form gives
\[
  e_p\equiv0\pmod2,
\]
so
\[
  n_p=e_p+1
\]
is odd.  Since $p$ is a $5$-input witness, the cyclotomic input condition gives
\[
  \ord_5(p)\mid n_p
\]
if $\ord_5(p)>1$, and gives
\[
  5\mid n_p
\]
if $\ord_5(p)=1$.

The possible orders modulo $5$ are
\[
  1,\ 2,\ 4.
\]
The orders $2$ and $4$ are even and therefore cannot divide the odd integer
$n_p$.  Hence the only possible non-Euler input order is
\[
  \ord_5(p)=1,
\]
which is equivalent to
\[
  p\equiv1\pmod5.
\]
In the order-one case the corrected cyclotomic input condition gives input
precisely when
\[
  5\mid e_p+1,
\]
and the valuation is
\[
  v_5(\sigma(p^{e_p}))=v_5(e_p+1).
\]

Now suppose that $p=\pi$ is the Euler prime, with exponent $\alpha$.  Put
\[
  n_\pi:=\alpha+1.
\]
Since
\[
  \alpha\equiv1\pmod4,
\]
we have
\[
  n_\pi\equiv2\pmod4.
\]

If
\[
  \ord_5(\pi)=1,
\]
then
\[
  \pi\equiv1\pmod5,
\]
and the order-one case gives input precisely when
\[
  5\mid \alpha+1.
\]
The valuation is then
\[
  v_5(\sigma(\pi^\alpha))=v_5(\alpha+1).
\]

If
\[
  \ord_5(\pi)=2,
\]
then
\[
  \pi\equiv-1\pmod5.
\]
Since
\[
  2\mid n_\pi,
\]
the cyclotomic input condition gives a $5$-adic contribution, and the valuation
formula is
\[
  v_5(\sigma(\pi^\alpha))
  =
  v_5(\pi^2-1)
  +
  v_5\!\left(\frac{\alpha+1}{2}\right).
\]

Finally, if
\[
  \ord_5(\pi)=4,
\]
then
\[
  \pi\equiv2\pmod5
  \quad\text{or}\quad
  \pi\equiv3\pmod5.
\]
But
\[
  n_\pi\equiv2\pmod4,
\]
so
\[
  4\nmid n_\pi.
\]
Thus no $5$-adic contribution occurs in these two residue classes.
\end{proof}

\end{document}
